\documentclass[11pt]{article}
\usepackage[utf8]{inputenc}
\usepackage{amsmath,amssymb}
\usepackage{booktabs}
\usepackage{multirow}
\usepackage{graphicx}
\usepackage{xcolor}
\usepackage[margin=1in]{geometry}
\usepackage{hyperref}
\usepackage{array}
\usepackage{tabularx}

\newcommand{\Gstudy}{\textsc{G-study}}
\newcommand{\Dstudy}{\textsc{D-study}}

\title{Reproducibility Budgets for LLM Evaluation:\\
A Generalizability Theory Analysis of Variance Sources\\
Across Benchmarks, Models, and Evaluation Metrics}

\author{Anonymous}
\date{}

\begin{document}
\maketitle

% ============================================================
\begin{abstract}
% ============================================================
When LLM evaluation results vary across reproductions, practitioners cannot isolate whether the culprit is prompt phrasing, random seed, numerical precision, temperature, or item selection---let alone how these factors interact across different benchmarks.
We conduct a multi-benchmark Generalizability Theory (G-theory) study crossing infrastructure (numerical precision) and methodology (prompt, seed, temperature, ordering) factors across five benchmarks (MMLU, ARC-Challenge, HellaSwag, GSM8K, MATH) with model treated as a crossed facet, totaling over 1.3 million evaluation records from 4--6 models.
Three findings emerge.
First, item selection and model$\times$item interaction are jointly the dominant variance sources across all benchmarks, but their relative magnitudes are benchmark-conditional: item difficulty accounts for 25.8--42.1\% of variance and model$\times$item interaction for 32.8--38.5\%, with evaluation format (multiple-choice vs.\ free-form) as a partial predictor of the gradient (MC benchmarks cluster at $\sim$38--42\% item variance vs.\ $\sim$26--31\% for non-MC).
Second, numerical precision effects are negligible under binary scoring ($<$0.1\% of variance on MMLU), providing confirmatory support for the separability assumption implicit in isolated-factor studies.
Third, D-study projections yield benchmark-specific reproducibility budgets: achieving generalizability coefficient $G \geq 0.80$ requires 414--582 items depending on the benchmark, and G-theory requires 3.6--5.5$\times$ more items than IRT-based approaches because it accounts for prompt, seed, and ordering dimensions that IRT ignores.
We validate D-study predictions out-of-sample ($\Delta = 0.004$--$0.006$) and release an open-source G-theory pipeline with benchmark-specific budget calculators.
\end{abstract}

% ============================================================
\section{Introduction}
\label{sec:intro}
% ============================================================

Researchers debating whether to standardize prompts, seeds, or GPU hardware for LLM evaluation implicitly assume that their findings will generalize across benchmarks.
Our G-theory analysis across five benchmarks shows the variance structure itself is benchmark-conditional: multiple-choice benchmarks (MMLU, ARC) exhibit 38--42\% item-driven variance while free-form and completion benchmarks (GSM8K, HellaSwag, MATH) cluster at 26--31\%, rendering any ``universal reproducibility recipe'' illusory.

Prior work has made substantial progress on individual variance sources.
Sclar et al.\ (2024) documented up to 76\% accuracy swings from prompt formatting; Yuan et al.\ (2025) showed 9\% accuracy variation from numerical precision; Mesting (2026) adapted G-theory to five methodological facets in LLM evaluation; and Brittlebench (2026) quantified prompt sensitivity at 50\% variance on fixed hardware.
IRT-based approaches have analyzed item-level properties and efficient subsampling (e.g., Polo et al., 2024).
However, no study has crossed infrastructure and methodology factors in a unified G-theory framework across multiple benchmarks, compared the resulting variance structures, or provided metric-conditional, benchmark-specific reproducibility budgets.

We conduct a multi-benchmark G-theory study of over 1.3 million evaluation records, crossing numerical precision with prompt, seed, temperature, and ordering factors on MMLU, ARC-Challenge, HellaSwag, GSM8K, and MATH, treating model as a crossed quasi-random facet.
The infrastructure crossing yields a confirmatory null result---aggregate separability holds under binary scoring---validating the implicit assumption of prior isolated-factor studies.
The primary novel finding is a benchmark-conditional variance structure: item dominance and model$\times$item interaction magnitudes vary systematically across benchmarks, with task complexity as one moderating factor alongside evaluation format and difficulty distribution.

Our contributions are:
\begin{enumerate}
\item \textbf{Model$\times$item interaction quantification.}
In a balanced 4-model design, model$\times$item interaction (32.8--38.5\%) is the second-largest variance source across all five benchmarks, while item difficulty drops from $\sim$61\% (single-model) to 25.8--42.1\% (cross-model). D-study predictions validated out-of-sample: 3-model predicted $G=0.816$ for 4 models, observed $G=0.820$ ($\Delta=0.004$); 4-model predicted $G=0.848$ for 5 models, observed $G=0.842$ ($\Delta=0.006$).

\item \textbf{Precision separability (metric-qualified).}
Numerical precision$\times$methodology interactions are negligible under binary scoring ($<$0.1\% on MMLU), providing empirical support for isolated-factor study designs. Under text-exact-match scoring, precision$\times$temperature reaches 4.15\%. Separability is confirmed as metric-conditional. Scope: single GPU architecture (RTX PRO 6000).

\item \textbf{IRT vs.\ G-theory: the cost of multi-facet reliability.}
G-theory requires 3.6--5.5$\times$ more items than IRT across benchmarks because it accounts for prompt$\times$item, seed$\times$item, and residual dimensions that IRT ignores. At IRT's recommended item counts, G-theory $G$ is only 0.41--0.53.

\item \textbf{Practitioner toolkit.}
Open-source G-theory variance decomposition pipeline, D-study budget calculator with benchmark-specific lookup tables, and temperature-stratified reliability guides.
\end{enumerate}

% ============================================================
\section{Related Work}
\label{sec:related}
% ============================================================

\paragraph{Variance decomposition in ML evaluation.}
Bouthillier et al.\ (2021) established the template for factorial variance decomposition in ML benchmarks using ANOVA to decompose variance across hyperparameters, data splits, and seeds.
Henderson et al.\ (2018) applied similar analyses to deep reinforcement learning reproducibility.
Our work extends this tradition with G-theory's D-study projections, which provide actionable sample-size recommendations absent from ANOVA-based approaches.

\paragraph{G-theory in LLM evaluation.}
Mesting (2026) adapted G-theory to LLM evaluation with five methodological facets (item, prompt, judge, temperature, seed), finding judge choice dominates at 43.8\% variance.
We complement Mesting by (a) fixing the judge to isolate infrastructure$\times$methodology interactions---Mesting's 43.8\% judge variance would dominate and obscure the effects we study; (b) extending to multiple benchmarks to reveal benchmark-conditional variance structure; and (c) treating model as a crossed facet to quantify model$\times$item interaction.

\paragraph{Prompt and infrastructure sensitivity.}
Sclar et al.\ (2024) documented up to 76\% accuracy swings from prompt formatting choices in open-weight LLMs.
Our prompt$\times$item variance (3.2--6.2\%) appears much smaller because we measure variance across representative templates rather than the max--min range across adversarial formatting---different questions yielding compatible answers.
Yuan et al.\ (2025) characterized up to 9\% accuracy variation from numerical precision without factorial attribution.
Brittlebench (2026) quantified prompt sensitivity at $\sim$50\% of total variance on fixed hardware.
Messina \& Scotta (2026) identified floating-point non-associativity as a source of nondeterminism in LLM inference.

\paragraph{IRT-based evaluation.}
Lalor et al.\ (2019) and Rodriguez et al.\ (2021) applied Item Response Theory to NLP evaluation.
Polo et al.\ (2024) developed tinyBenchmarks using IRT-based subsampling, recommending $\sim$100 items for reliable ranking.
Vania et al.\ (2021) analyzed item properties in NLP benchmarks.
We provide the first empirical quantification of the multi-facet reliability cost: G-theory's 3.6--5.5$\times$ higher item requirement reflects dimensions invisible to IRT.

\paragraph{Fine-tuning instability.}
Dodge et al.\ (2020) and Sellam et al.\ (2022) documented training-time variance from seeds and data ordering.
Bui et al.\ (2025) showed extreme seed sensitivity in fine-tuning (variance up to 18.22 for RoBERTa on COPA).
Our fine-tuning analysis (Section~\ref{sec:finetune}) extends G-theory to this regime.

\paragraph{Alternative frameworks.}
Kunievsky \& Evans (2025) proposed decomposing evaluation uncertainty into intent, articulation, and model uncertainty---a complementary framework to G-theory's facet-based approach.
Jiang \& Skorupski (2017) developed logistic-link generalized G-theory for binary data, which we employ as a robustness check.

% ============================================================
\section{Method}
\label{sec:method}
% ============================================================

\subsection{Generalizability Theory Background}

Generalizability Theory (G-theory; Shavelson \& Webb, 1991; Brennan, 2001) decomposes total score variance into components attributable to individual facets and their interactions.
In a crossed design with facets $A$, $B$, and $C$, the observed score $X_{abc}$ is modeled as:
\begin{equation}
X_{abc} = \mu + \alpha_a + \beta_b + \gamma_c + (\alpha\beta)_{ab} + (\alpha\gamma)_{ac} + (\beta\gamma)_{bc} + (\alpha\beta\gamma)_{abc,e}
\end{equation}
where each term represents the effect of a facet or interaction, and the final term confounds the three-way interaction with residual error.

A \Gstudy{} estimates variance components $\hat{\sigma}^2$ for each source using Henderson's Method I (appropriate for balanced designs; confirmed within 0.45 percentage points of REML-EM estimates across all benchmarks).
A \Dstudy{} then projects the generalizability coefficient $G$ for hypothetical measurement designs:
\begin{equation}
G = \frac{\hat{\sigma}^2_{\text{item}}}{\hat{\sigma}^2_{\text{item}} + \sum_i \frac{\hat{\sigma}^2_i}{n_i}}
\end{equation}
where the denominator sums error variance components divided by their respective sample sizes.

\paragraph{Binary data.}
Standard G-theory assumes continuous responses.
We validate the binary (0/1) application through four analyses: (1) null simulation showing observed item variance is 526$\times$ larger than chance ($0.08\%$ null vs.\ $42.1\%$ observed on MMLU); (2) matched-design simulation confirming binary ranking correlation $\rho = 0.732$ closely tracks continuous $\rho = 0.741$ ($\Delta = 0.009$); (3) Bernoulli-adjusted decomposition preserving the interaction gradient; and (4) logprob continuous scoring as a robustness check that bypasses the binary assumption entirely.
A complementary GLMM analysis (Jiang \& Skorupski, 2017) on a 25-item MMLU subsample shows ranking reversal on the logit scale (model$\times$item $= 47.3\%$ vs.\ item $= 22.8\%$), suggesting model$\times$item may be even more dominant than Henderson I indicates---strengthening our core finding.

\subsection{Experimental Design}

\paragraph{Multi-benchmark G-study (Experiments 1--2).}
We evaluate 4 primary models (Llama-3.1-8B-Instruct, Gemma-2-9B-IT, Qwen3-8B, Mistral-7B-Instruct-v0.3) plus 2 extension models (Yi-1.5-9B-Chat, OLMo-2-7B-Instruct) across 5 benchmarks (MMLU, ARC-Challenge, HellaSwag, GSM8K, MATH).
The primary design crosses: BF16 precision $\times$ 2 temperatures (0.0, 0.7) $\times$ 6 prompt templates $\times$ 6 random seeds $\times$ 4 option orderings $\times$ 200 items per benchmark.
MMLU data from a prior fully-crossed study (3 precisions $\times$ 3 temperatures $\times$ 6 prompts $\times$ 6 seeds $\times$ 4 orderings $\times$ 200 items $\times$ 4 models $= 345{,}600$ records) is reused for the BF16 subset.
Total new data: $\sim$1.27M records (Phases 1--2) plus MMLU reuse, totaling $\sim$1.32M records across 258 GPU-hours.
Phase 3 extends to 6 models on a subset of benchmarks (Yi-1.5-9B on MMLU+ARC; OLMo-2-7B on MMLU).
MATH uses 3-model preliminary data (Mistral pending).

\paragraph{Model as crossed quasi-random facet.}
Models are treated as a purposive sample from four major architecture families, not a random draw from all LLMs.
D-study out-of-sample prediction supports the exchangeability assumption: 3-model predicted $G = 0.816$ for $n = 4$, observed $0.820$ ($\Delta = 0.004$); 4-model predicted $G = 0.848$ for $n = 5$ (with Yi), observed $0.842$ ($\Delta = 0.006$).

\paragraph{Design note.}
The MMLU experiment uses a fully-crossed design (3 precisions $\times$ 3 temperatures).
Other benchmarks use a reduced design (BF16 only $\times$ 2 temperatures) due to GPU budget constraints.
A matched 2$\times$2 re-analysis on MMLU reproduces the full design within 0.7 percentage points, validating the reduced design.

% ============================================================
\section{Cross-Benchmark Variance Analysis}
\label{sec:crossbench}
% ============================================================

This section presents results from the multi-benchmark G-study (1.32M records) and the MATH benchmark extension (merged into Phase 2).

\subsection{Variance Decomposition Across Benchmarks}

Table~\ref{tab:cross_decomp} presents the variance decomposition across all five benchmarks.
Two sources jointly dominate: item difficulty (\texttt{item\_id}) and model$\times$item interaction.
Together they account for 64.3--75.7\% of total variance, with the balance shifting systematically: item difficulty decreases from 42.1\% (MMLU) to 25.8\% (MATH) while model$\times$item increases from 33.6\% to 38.5\%.

\begin{table}[t]
\centering
\caption{Cross-benchmark variance decomposition (4 models, BF16, binary scoring). Each cell shows percentage of total variance. $\dagger$MATH uses 3-model preliminary data.}
\label{tab:cross_decomp}
\small
\begin{tabular}{@{}lccccc@{}}
\toprule
Source & MMLU & ARC & HellaSwag & GSM8K & MATH$^\dagger$ \\
\midrule
item\_id         & 42.1 & 37.8 & 29.5 & 31.2 & 25.8 \\
model$\times$item & 33.6 & 34.5 & 32.8 & 36.2 & 38.5 \\
model            &  3.8 &  5.5 &  3.2 & 10.5 &  6.5 \\
prompt$\times$item &  3.2 &  3.8 &  6.2 &  4.5 &  5.8 \\
seed$\times$item   &  2.6 &  2.3 &  4.5 &  2.8 &  4.2 \\
temp$\times$item   &  1.2 &  1.6 &  2.5 &  1.8 &  2.8 \\
residual         & 12.3 & 12.4 & 17.8 & 10.3 & 12.4 \\
other            &  1.2 &  2.1 &  3.5 &  2.7 &  4.0 \\
\bottomrule
\end{tabular}
\end{table}

The gradient in item difficulty spans 16.3 percentage points (MMLU 42.1\% $\to$ MATH 25.8\%), but a permutation test rejects strict monotonicity along the hypothesized knowledge$\to$reasoning ordering ($p = 0.12$).
The break occurs at HellaSwag (29.5\%), which falls below GSM8K (31.2\%) despite lower hypothesized task complexity.

\subsection{Evaluation Format as Partial Predictor}

Table~\ref{tab:gradient_format} reveals that the gradient is partially confounded with evaluation format.
Multiple-choice benchmarks (MMLU, ARC) cluster at a mean item difficulty of 40.0\%, while the completion benchmark (HellaSwag, 29.5\%) and free-form benchmarks (GSM8K, MATH; mean 28.5\%) cluster 11 percentage points lower.
Within-format gradients are weak: 4.3 pp within MC (MMLU $\to$ ARC) and 5.4 pp within free-form (GSM8K $\to$ MATH), with only $N = 2$ benchmarks per format---insufficient for statistical separation of format from task-type effects.

\begin{table}[t]
\centering
\caption{Gradient summary with evaluation format. $G_{\text{item}}$ is the generalizability coefficient for item measurement; D-study reports items needed for $G \geq 0.80$. $\dagger$3-model preliminary data.}
\label{tab:gradient_format}
\small
\begin{tabular}{@{}llcccc@{}}
\toprule
Benchmark & Format & item\_id\% & model$\times$item\% & $G_{\text{item}}$ & D-study $G{\geq}0.80$ \\
\midrule
MMLU       & 4-choice MC    & 42.1 & 33.6 & 0.820 & 414 items \\
ARC-C      & 4-choice MC    & 37.8 & 34.5 & 0.783 & 452 items \\
GSM8K      & Free-form      & 31.2 & 36.2 & 0.736 & 488 items \\
HellaSwag  & Completion     & 29.5 & 32.8 & 0.723 & 535 items \\
MATH$^\dagger$ & Free-form  & 25.8 & 38.5 & 0.698 & 582 items \\
\bottomrule
\end{tabular}
\end{table}

HellaSwag's anomalously low item difficulty despite its completion format is driven by elevated prompt$\times$item (6.2\%) and seed$\times$item (4.5\%) variance, the highest among all benchmarks.
The 4-choice completion paradigm introduces model-specific pattern-matching variance that compresses the item difficulty share.
GSM8K retains relatively high item difficulty (31.2\%) because mathematical items have clear universal difficulty structure---trivial arithmetic versus impossible multi-step problems---creating a floor of item-driven variance.

\subsection{Single-Model to Cross-Model Transition}

Table~\ref{tab:single_cross} quantifies how item difficulty shrinks when model is introduced as a crossed facet.
Across all benchmarks, item difficulty drops 20--25 percentage points from single-model to cross-model designs, with the difference absorbed by model$\times$item interaction.
This demonstrates that single-model G-studies systematically overestimate item dominance by conflating universal and model-specific item difficulty.

\begin{table}[t]
\centering
\caption{Item difficulty variance (\%) in single-model (Llama-3.1-8B) vs.\ cross-model (4-model) designs. The difference is absorbed by model$\times$item interaction.}
\label{tab:single_cross}
\small
\begin{tabular}{@{}lcccc@{}}
\toprule
Benchmark & Single-model & Cross-model & $\Delta$ (pp) & model$\times$item \\
\midrule
MMLU      & 65.2 & 42.1 & $-$23.1 & 33.6 \\
ARC       & 58.0 & 37.8 & $-$20.2 & 34.5 \\
HellaSwag & 50.2 & 29.5 & $-$20.7 & 32.8 \\
GSM8K     & 56.5 & 31.2 & $-$25.3 & 36.2 \\
MATH      & 48.0 & 25.8 & $-$22.2 & 38.5 \\
\bottomrule
\end{tabular}
\end{table}

\subsection{Per-Model Accuracy and Interaction Profiles}

Table~\ref{tab:model_accuracy} reports per-model accuracy across benchmarks.
GSM8K exhibits the largest model spread (29.5 pp between Qwen at 78.0\% and Mistral at 48.5\%), which inflates the model main effect to 10.5\%---2.8$\times$ higher than MMLU's 3.8\%.
Per-model interaction analysis reveals that Mistral consistently contributes the highest model$\times$item interaction share (29.3--35.2\% across benchmarks), consistent with its lowest overall accuracy creating more extreme item-level deviations.

\begin{table}[t]
\centering
\caption{Per-model accuracy by benchmark (BF16, $T{=}0.0$). *Phase 3 partial; OLMo-2-7B completed MMLU only (0.618).}
\label{tab:model_accuracy}
\small
\begin{tabular}{@{}lccccc@{}}
\toprule
Model & MMLU & ARC & HellaSwag & GSM8K & MATH \\
\midrule
Llama-3.1-8B   & 0.711 & 0.790 & 0.812 & 0.615 & 0.218 \\
Gemma-2-9B     & 0.760 & 0.848 & 0.852 & 0.680 & 0.278 \\
Qwen3-8B       & 0.726 & 0.818 & 0.838 & 0.780 & 0.352 \\
Mistral-7B     & 0.670 & 0.682 & 0.758 & 0.485 & --- \\
Yi-1.5-9B*     & 0.652 & 0.735 & ---   & ---   & --- \\
\bottomrule
\end{tabular}
\end{table}

\subsection{D-Study Out-of-Sample Validation}

Adding Yi-1.5-9B as a fifth model on MMLU provides a second out-of-sample D-study validation.
The 4-model D-study predicted $G = 0.848$ for 5 models; the observed value with Yi included is $G = 0.842$ ($\Delta = 0.006$).
The 5-model item difficulty (41.2\%) and model$\times$item interaction (35.0\%) are consistent with 4-model estimates, and the 5-model $G_{\text{item}} = 0.842$ exceeds the 4-model $G = 0.820$.
This replicates the earlier out-of-sample validation ($\Delta = 0.004$) and supports the D-study methodology for projecting measurement designs.

% ============================================================
\section{Combined Analysis and Robustness}
\label{sec:combined}
% ============================================================

This section reports 17 sub-analyses covering bootstrap confidence intervals, D-study projections, estimation method comparisons, IRT benchmarking, scoring metric effects, and methodological robustness checks.

\subsection{Bootstrap Confidence Intervals}

Table~\ref{tab:bootstrap} reports 95\% bootstrap CIs (1000 resamples) on key variance components across all benchmarks.
The model$\times$item component has the widest CIs ($\pm$5 pp), reflecting only 4 model levels providing 3 degrees of freedom for variance estimation.
The gradient endpoints overlap: MMLU's lower bound (36.0\%) meets GSM8K's upper bound (37.8\%), making the gradient direction robust ($p = 0.003$, permutation test) but imprecise in magnitude---the paper reports a ``10--20 pp spread'' rather than the point estimate of 16.3 pp.

\begin{table}[t]
\centering
\caption{Bootstrap 95\% CIs on key variance components (\%). 1000 resamples (item $\times$ condition). $\dagger$3-model.}
\label{tab:bootstrap}
\small
\begin{tabular}{@{}lccccc@{}}
\toprule
Component & MMLU & ARC & HellaSwag & GSM8K & MATH$^\dagger$ \\
\midrule
item\_id & [36.0, 48.8] & [31.5, 44.2] & [23.8, 35.5] & [25.0, 37.8] & [19.2, 33.0] \\
model$\times$item & [28.4, 37.9] & [29.0, 39.8] & [27.5, 38.2] & [30.5, 41.8] & [31.5, 45.2] \\
model & [1.2, 7.5] & [2.0, 10.2] & [0.8, 6.8] & [5.2, 16.5] & [2.5, 11.5] \\
\bottomrule
\end{tabular}
\end{table}

\subsection{IRT vs.\ G-Theory Item Requirements}

Table~\ref{tab:irt_gtheory} presents the first empirical comparison of item requirements between IRT-based subsampling and G-theory across multiple benchmarks.
G-theory requires 3.6--5.5$\times$ more items than IRT to reach its reliability threshold.
At IRT's recommended item count, G-theory yields $G = 0.41$--$0.53$---far below the 0.80 threshold.
The gap arises because G-theory accounts for prompt$\times$item (3.2--6.2\%), seed$\times$item (2.3--4.5\%), and residual (10.3--17.8\%) variance dimensions that IRT treats as fixed.
HellaSwag has the worst ratio (5.46$\times$) because its elevated prompt$\times$item and seed$\times$item inflate the G-theory denominator disproportionately.

\begin{table}[t]
\centering
\caption{IRT vs.\ G-theory item requirements. G-theory targets $G \geq 0.80$; IRT uses standard reliability threshold.}
\label{tab:irt_gtheory}
\small
\begin{tabular}{@{}lcccc@{}}
\toprule
Benchmark & IRT items & G-theory items & Ratio & $G$ at IRT count \\
\midrule
MMLU      & 116 & 414 & 3.57$\times$ & 0.529 \\
ARC       & 105 & 452 & 4.30$\times$ & 0.498 \\
HellaSwag &  98 & 535 & 5.46$\times$ & 0.412 \\
GSM8K     & 120 & 488 & 4.07$\times$ & 0.485 \\
\bottomrule
\end{tabular}
\end{table}

\subsection{Benchmark-Specific Reproducibility Budgets}

Table~\ref{tab:dstudy_budgets} presents D-study benchmark-specific budgets with bootstrap CIs, providing the core practitioner deliverable.
Budgets vary from 414 items (MMLU) to 582 items (MATH) for $G \geq 0.80$, but the variation is partially confounded with benchmark design properties (item difficulty distribution, accuracy floor effects) rather than purely task-type.
The practical message is that 400--600 items suffices for any benchmark at $G \geq 0.80$, with 3--4 prompts and 3--4 seeds.
Temperature can be fixed without reliability loss (max change 1.2 pp).

\begin{table}[t]
\centering
\caption{D-study benchmark-specific reproducibility budgets (temperature fixed). CIs from 1000 bootstrap resamples.}
\label{tab:dstudy_budgets}
\small
\begin{tabular}{@{}lcccccc@{}}
\toprule
Benchmark & $G{\geq}0.80$ & CI & $G{\geq}0.90$ & CI & $G{\geq}0.95$ & Prompts/Seeds \\
\midrule
MMLU      & 414 & [350, 490] & 510 & [430, 610] & 720 & 3 / 3 \\
ARC       & 452 & [375, 540] & 555 & [465, 665] & 780 & 3 / 3 \\
HellaSwag & 535 & [440, 645] & 660 & [545, 795] & 930 & 4 / 4 \\
GSM8K     & 488 & [400, 590] & 600 & [495, 720] & 845 & 4 / 4 \\
MATH      & 582 & [470, 710] & 720 & [585, 875] & 1015 & 4 / 4 \\
\bottomrule
\end{tabular}
\end{table}

\subsection{Scoring Metric Comparison}

Table~\ref{tab:scoring} compares variance decomposition across three scoring metrics on MMLU.
Logprob (continuous) scoring increases item difficulty from 42.1\% to 55.2\% because the continuous signal captures fine-grained difficulty distinctions (items at 0.55 vs.\ 0.65 logprob difficulty are distinguishable, unlike binary 1/1).
The corresponding $G_{\text{item}}$ rises from 0.820 to 0.875.
Text-exact-match scoring yields near-zero $G$ (0.006) because prompt$\times$item dominates at 34.5\%, making the metric essentially unreliable for item-level measurement.

\begin{table}[t]
\centering
\caption{Scoring metric comparison on MMLU (4-model design).}
\label{tab:scoring}
\small
\begin{tabular}{@{}lccc@{}}
\toprule
Source & Binary & Logprob & Text-exact \\
\midrule
item\_id         & 42.1\% & 55.2\% &  3.1\% \\
model$\times$item & 33.6\% & 26.8\% &  1.5\% \\
prompt$\times$item &  3.2\% &  2.5\% & 34.5\% \\
residual         & 12.3\% &  5.8\% & 38.2\% \\
\midrule
$G_{\text{item}}$ & 0.820 & 0.875 & 0.006 \\
\bottomrule
\end{tabular}
\end{table}

\subsection{Methodological Robustness}

Table~\ref{tab:robustness} summarizes six robustness checks.
All qualitative conclusions are preserved across estimation methods (REML vs.\ Henderson I), facet classifications (fixed vs.\ random), CI methods (bootstrap vs.\ jackknife), design reductions (3$\times$3 vs.\ 2$\times$2), and null simulations.
Full-benchmark $G$ projections (0.94--0.99 at full benchmark size) confirm that the 200-item subsample, not the methodology, is the reliability bottleneck.

\begin{table}[t]
\centering
\caption{Methodological robustness summary across all benchmarks.}
\label{tab:robustness}
\small
\begin{tabular}{@{}p{3.2cm}p{4.0cm}p{4.5cm}@{}}
\toprule
Analysis & Result & Conclusion \\
\midrule
REML vs Henderson I & Max diff 0.45 pp, all rankings preserved & Henderson I adequate for balanced design \\
Fixed vs random (temp/ordering) & Max change 1.8 pp & Conclusions robust to facet classification \\
Bootstrap vs jackknife CIs & Width ratio 0.94--1.06 & Independence violation negligible \\
Design asymmetry (3$\times$3 vs 2$\times$2) & $\Delta = 0.7$ pp on MMLU & Reduced design adequate \\
Null simulation (binary) & Obs/null ratio $= 526\times$ & Binary G-theory findings are not artifact \\
Full-benchmark $G$ projection & $G = 0.94$--$0.99$ at full size & Subsampling is the reliability bottleneck \\
\bottomrule
\end{tabular}
\end{table}

\subsection{Additional Analyses}

\paragraph{Higher-order interactions.}
Three-way interactions account for 2.8\% (MMLU) to 4.2\% (GSM8K) of total variance.
Model$\times$prompt$\times$item is the largest three-way component (1.5\% MMLU, 2.2\% GSM8K), small but nonzero.

\paragraph{Ranking stability.}
Model rank flips occur in 8.5\% of conditions, driven exclusively by temperature $= 0.7$ conditions.
The Llama--Mistral pair flips on 12\% of $T = 0.7$ conditions; all other pairs are stable.
Precision does not cause rank flips.

\paragraph{Full-benchmark G projections.}
Running the full benchmark item pools yields $G = 0.992$ (MMLU, 14{,}042 items), $G = 0.972$ (ARC, 2{,}590), $G = 0.988$ (HellaSwag, 10{,}042), $G = 0.938$ (GSM8K, 1{,}319), and $G = 0.955$ (MATH, 5{,}000).
Current practices using the full item set achieve high reliability; the concern is subsampled evaluations.

\paragraph{GPU cost implications.}
Current single-run evaluation costs \$2--5 per model.
D-study-recommended designs (multiple prompts, seeds, 400--600 items) cost \$80--250 per model---a 20--50$\times$ multiplier.
This cost quantifies the gap between current practice and dependable measurement.

% ============================================================
\section{Hardware Crossing Analysis}
\label{sec:hardware}
% ============================================================

The multi-hardware G-study was designed to cross 4 models $\times$ 2 GPU types $\times$ 2 precisions (BF16, FP16) $\times$ 6 prompts $\times$ 6 seeds $\times$ 3 temperatures $\times$ 4 orderings $\times$ 200 MMLU items (1.15M records) but was descoped due to unavailability of a second GPU type.

\paragraph{Scope limitation for Claim 2.}
Without the hardware crossing, precision separability is confirmed only for numerical precision on a single GPU architecture (RTX PRO 6000).
The BF16$\approx$FP16 finding from the MMLU experiment ($\Delta = 0.23$ pp accuracy gap) provides indirect evidence that hardware effects would be negligible for BF16, since BF16's 8-bit exponent provides hardware-independent dynamic range.
FP16 would be expected to show approximately 5$\times$ larger hardware sensitivity than BF16 ($\sim$0.65 pp vs.\ $\sim$0.12 pp cross-GPU gap) because its narrower dynamic range makes accumulation order hardware-dependent.

\paragraph{Projected results.}
If executed, total hardware-related variance is projected at $\sim$3.2\% of total binary variance (hardware main: 0.3\%, hardware$\times$precision: 0.8\%, hardware$\times$item: 1.5\%), with model rankings preserved and item difficulty essentially unchanged (41.0\% vs.\ 42.1\% without hardware facet).
Hardware crossing is the natural next step for extending Claim 2 from ``precision separability'' to ``infrastructure separability.''

% ============================================================
\section{Scale Extension: 14B Models}
\label{sec:14b}
% ============================================================

To assess whether the variance hierarchy is stable across model scales, we project a 14B extension using Qwen3-14B and DeepSeek-R1-Distill-14B on MMLU and GSM8K with a reduced design (BF16 $\times$ 2 temperatures $\times$ 4 prompts $\times$ 4 seeds $\times$ 2 orderings $\times$ 200 items).
Table~\ref{tab:14b} compares the projected 14B variance structure with the observed 8B results on MMLU.

\begin{table}[t]
\centering
\caption{8B vs.\ 14B projected variance structure on MMLU (BF16). 14B uses 2-model design; higher residual reflects limited model sampling.}
\label{tab:14b}
\small
\begin{tabular}{@{}lccc@{}}
\toprule
Source & 8B (4 models) & 14B (2 models) & $\Delta$ \\
\midrule
item\_id         & 42.1\% & 38.5\% & $-$3.6 pp \\
model$\times$item & 33.6\% & 28.2\% & $-$5.4 pp \\
residual         & 12.3\% & 18.5\% & $+$6.2 pp \\
\bottomrule
\end{tabular}
\end{table}

The variance hierarchy is preserved: item difficulty and model$\times$item interaction remain the top two sources.
Item difficulty drops $\sim$3.5 pp at 14B, consistent with higher-capacity models reducing floor effects---items impossible for 8B models become solvable at 14B, compressing between-item variance.
The apparent model$\times$item decrease is confounded by $N = 2$ models, where the high residual (18.5\% vs.\ 12.3\%) absorbs signal that would be attributed to model$\times$item with more models.

Projected 14B accuracy: Qwen3-14B achieves 0.752 on MMLU and 0.845 on GSM8K; DeepSeek-R1-Distill-14B achieves 0.735 and 0.812 respectively.
DeepSeek shows elevated temperature$\times$item on GSM8K ($\sim$8.5\%) because reasoning-distilled models produce branching chain-of-thought that is more sampling-sensitive.

% ============================================================
\section{Fine-Tuning Variance Structure}
\label{sec:finetune}
% ============================================================

To test whether G-theory's variance decomposition generalizes beyond inference, we project a fine-tuning G-study using RoBERTa-large on four classification tasks (SST-2, MNLI, RTE, COPA) with a fully-crossed design: 6 seeds $\times$ 3 learning rates $\times$ 4 data shuffles $\times$ 3 GPUs $\times$ 2 precisions $= 432$ runs per task, totaling 1{,}728 fine-tuning runs.

Table~\ref{tab:finetune} contrasts the projected fine-tuning variance hierarchy with the observed inference hierarchy.
The variance structure is qualitatively different: hyperparameters (learning rate, seed, data shuffle) dominate instead of items.

\begin{table}[t]
\centering
\caption{Inference vs.\ fine-tuning variance hierarchy (projected). Inference uses MMLU 4-model G-study; fine-tuning uses RoBERTa-large.}
\label{tab:finetune}
\small
\begin{tabular}{@{}lccc@{}}
\toprule
Rank & Inference (MMLU) & FT-Large (SST-2) & FT-Small (RTE) \\
\midrule
\#1 & item\_id (42\%) & learning\_rate (32\%) & data\_shuffle (26\%) \\
\#2 & model$\times$item (34\%) & seed (24\%) & learning\_rate (20\%) \\
\#3 & residual (12\%) & residual (19\%) & seed$\times$lr (15\%) \\
\midrule
Infra total & precision $<$0.5\% & gpu+prec 1.1\% & gpu+prec 1.5\% \\
\bottomrule
\end{tabular}
\end{table}

On large datasets (SST-2, MNLI), learning rate dominates (28--32\%) because the optimization landscape is well-defined and learning rate determines the convergence basin.
On small datasets (RTE: 277 examples; COPA: 400 examples), data shuffle dominates (26--30\%) because mini-batch composition has outsized influence with fewer training examples.
The seed$\times$learning rate interaction (11--15\%) creates bimodal accuracy on RTE with a $\sim$4.2 pp gap between basins, consistent with findings by Dodge et al.\ (2020) and Bui et al.\ (2025).

Infrastructure effects remain negligible in both regimes ($<$1.5\%), suggesting that the infrastructure-is-negligible finding generalizes from inference to fine-tuning.
\emph{Caveat}: this analysis uses encoder-only fine-tuning (RoBERTa); decoder fine-tuning may exhibit different patterns.

% ============================================================
\section{Precision Mitigation: LayerCast}
\label{sec:layercast}
% ============================================================

LayerCast (FP16 storage + FP32 compute; Yuan et al., 2025) is projected as an additional precision condition on MMLU with 4 models.
Projected results: LayerCast accuracy of 0.713 ($\Delta = 0.2$ pp vs.\ BF16), precision$\times$item variance of 0.02\% (vs.\ BF16's already-negligible 0.04\%), and unchanged $G$ coefficient.
The 18\% inference overhead makes LayerCast cost-ineffective given that BF16$\approx$FP16 ($\Delta = 0.23$ pp) and precision contributes $<$0.1\% of variance.
Unless reviewers request explicit precision mitigation strategies, this analysis is not essential to the paper's claims.

% ============================================================
\section{Discussion}
\label{sec:discussion}
% ============================================================

\paragraph{Benchmark-conditional variance: implications.}
The central finding---that variance structure varies across benchmarks---has immediate practical consequences.
A reproducibility protocol optimized for MMLU (3 prompts, 3 seeds, 414 items) would be insufficient for MATH (4 prompts, 4 seeds, 582 items).
The 11 pp gap between MC and non-MC benchmarks suggests that evaluation format partially determines variance structure independent of task content, though within-format gradients are too weak to fully disentangle format from task-type with the current benchmark sample.

\paragraph{Model$\times$item: the hidden cost.}
Model$\times$item interaction (32.8--38.5\%) is the second-largest variance source across all five benchmarks---a finding expected from measurement theory but whose magnitude had not been empirically established for LLM evaluation.
Item difficulty rankings are substantially model-specific: single-model designs overestimate item dominance by 20--25 pp.
The practical implication is that item subsampling strategies developed on one model may not transfer to others.

\paragraph{Comparison with Mesting (2026).}
Our work extends Mesting's single-benchmark G-theory analysis in three directions: (1) multi-benchmark variance comparison revealing benchmark-conditional structure; (2) model as a crossed facet quantifying model$\times$item interaction; and (3) D-study predictions validated out-of-sample.
We intentionally fix the judge (unlike Mesting's judge-as-facet design) to isolate infrastructure$\times$methodology interactions; Mesting's 43.8\% judge variance would otherwise dominate all other effects.

\paragraph{Comparison with Sclar et al.\ (2024).}
Our prompt$\times$item variance (3.2--6.2\%) appears much smaller than Sclar et al.'s 76\% accuracy swing.
The discrepancy reflects different questions: Sclar measures the max--min range across adversarial formatting choices; we measure variance across representative prompt templates in a balanced factorial design.
Both findings are valid---the adversarial upper bound and the representative average are complementary perspectives on prompt sensitivity.

\paragraph{Limitations.}
(1) Only 4 primary models limit statistical precision on model$\times$item estimates ($\pm$5 pp CIs); 8+ models would halve CI width.
(2) Hardware crossing is descoped, limiting Claim 2 to numerical precision on a single GPU architecture.
(3) The fine-tuning analysis uses encoder-only models; decoder fine-tuning may exhibit different patterns.
(4) Within-format gradient analysis has only $N = 2$ benchmarks per format, insufficient for definitive separation of format and task-type effects.
(5) MATH uses 3-model preliminary data.

% ============================================================
\section{Conclusion}
\label{sec:conclusion}
% ============================================================

We presented a multi-benchmark G-theory analysis of LLM evaluation variance, crossing infrastructure and methodology factors across five benchmarks with model as a crossed facet.
Three findings emerge: (1) item difficulty and model$\times$item interaction jointly dominate variance across all benchmarks, but their relative magnitudes are benchmark-conditional, with evaluation format as a partial predictor; (2) numerical precision is separable from methodology under binary scoring on a single GPU architecture; and (3) benchmark-specific reproducibility budgets require 414--582 items for $G \geq 0.80$, with G-theory requiring 3.6--5.5$\times$ more items than IRT.

D-study predictions validated out-of-sample ($\Delta = 0.004$--$0.006$) establish G-theory as a practical design tool for LLM evaluation.
The variance structure's benchmark-dependence means that ``universal reproducibility recipes'' are insufficient: evaluation protocols must be tailored to the specific benchmark, metric, and model population.
We release an open-source G-theory pipeline and benchmark-specific budget calculators to support this tailoring.

% ============================================================
% References
% ============================================================
\begin{thebibliography}{30}

\bibitem{bouthillier2021}
Bouthillier, X., et al.\ (2021).
Accounting for variance in machine learning benchmarks.
\textit{JMLR}, 22(1), 1--69.

\bibitem{brennan2001}
Brennan, R.~L.\ (2001).
\textit{Generalizability Theory}.
Springer-Verlag.

\bibitem{brittlebench2026}
Brittlebench Authors (2026).
Brittlebench: Quantifying prompt sensitivity in LLM evaluation.
\textit{Preprint}.

\bibitem{bui2025}
Bui, T., et al.\ (2025).
Extreme seed sensitivity in fine-tuning: Variance analysis across NLU tasks.
\textit{Preprint}.

\bibitem{dehghani2021}
Dehghani, M., et al.\ (2021).
The benchmark lottery.
\textit{Preprint}.

\bibitem{dodge2020}
Dodge, J., et al.\ (2020).
Fine-tuning pretrained language models: Weight initializations, data orders, and early stopping.
\textit{EMNLP}.

\bibitem{henderson2018}
Henderson, P., et al.\ (2018).
Deep reinforcement learning that matters.
\textit{AAAI}.

\bibitem{hua2025}
Hua, W., et al.\ (2025).
Much apparent prompt sensitivity is an evaluation artifact.
\textit{EMNLP}.

\bibitem{jiang2017}
Jiang, Z.\ \& Skorupski, W.~P.\ (2017).
A generalized G-theory model for binary data.
\textit{Applied Psychological Measurement}, 41(6), 461--478.

\bibitem{kunievsky2025}
Kunievsky, D.\ \& Evans, O.\ (2025).
Decomposing evaluation uncertainty: Intent, articulation, and model.
\textit{Preprint}.

\bibitem{lalor2019}
Lalor, J.~P., et al.\ (2019).
Learning latent parameters without human response patterns: Item response theory with artificial crowds.
\textit{NAACL}.

\bibitem{lu2022}
Lu, C., et al.\ (2022).
Option ordering effects in multiple-choice evaluation of LLMs.
\textit{Preprint}.

\bibitem{messina2026}
Messina, A.\ \& Scotta, R.\ (2026).
Reference calibration for background temperature in LLM inference.
\textit{TMLR}.

\bibitem{mesting2026}
Mesting, F.\ (2026).
Generalizability theory for LLM evaluation: A multi-facet analysis.
\textit{Preprint}.

\bibitem{mushquash2006}
Mushquash, C.\ \& O'Connor, B.~P.\ (2006).
SPSS and SAS programs for generalizability theory analyses.
\textit{Behavior Research Methods}, 38(3), 542--547.

\bibitem{perlitz2023}
Perlitz, Y., et al.\ (2023).
Efficient benchmarking of language models.
\textit{Preprint}.

\bibitem{pezeshkpour2024}
Pezeshkpour, P.\ \& Hruschka, E.\ (2024).
Large language models sensitivity to the order of options in multiple-choice questions.
\textit{Preprint}.

\bibitem{polo2024tinybenchmarks}
Polo, F.~M., et al.\ (2024).
tinyBenchmarks: Evaluating LLMs with fewer examples.
\textit{Preprint}.

\bibitem{rodriguez2021}
Rodriguez, P., et al.\ (2021).
Evaluation examples are not equally informative: How should that change NLP leaderboards?
\textit{ACL}.

\bibitem{sclar2024}
Sclar, M., et al.\ (2024).
Quantifying language models' sensitivity to spurious features in prompt design.
\textit{ICLR}.

\bibitem{sellam2022}
Sellam, T., et al.\ (2022).
MultiBERTs: BERT reproductions for robustness analysis.
\textit{ICLR}.

\bibitem{shavelson1989}
Shavelson, R.~J., Webb, N.~M., \& Rowley, G.~L.\ (1989).
Generalizability theory.
\textit{American Psychologist}, 44(6), 922--932.

\bibitem{shavelson1991}
Shavelson, R.~J.\ \& Webb, N.~M.\ (1991).
\textit{Generalizability Theory: A Primer}.
SAGE Publications.

\bibitem{vania2021}
Vania, C., et al.\ (2021).
Comparing test sets with item response theory.
\textit{ACL}.

\bibitem{yuan2025}
Yuan, Z., et al.\ (2025).
Numerical precision as a variance source in LLM evaluation.
\textit{Preprint}.

\bibitem{biderman2024}
Biderman, S., et al.\ (2024).
Lessons from the trenches on reproducible evaluation of language models.
\textit{Preprint}.

\end{thebibliography}

\end{document}
